\begin{pastexam}[label={px:ntut-111-ci-02}]{北科大資工所 111 學年度 第三題}
（是非）

(10 pts) Mark by T(=true) or F(=False) each of the following statements. You don't need to prove it.
\begin{enumerate}[label=(\arabic*)]
\item $3n^3 + 7n^2 = \Omega(n^2)$ (2 pts)
\item $n^3 / \log^2 n = \Theta(n^3)$ (2 pts)
\item $f(n) = \mathrm{O}((f(n))^2)$ (2 pts)
\item When deleting a key from an AVL tree, at most one restructuring operation can restore the height balance property of the binary tree. (2 pts)
\item The operation insert() on a priority queue realized by a sorted list takes linear time. (2 pts)
\end{enumerate}
\end{pastexam}

\begin{pastsoln}
答案：(1) T；(2) F；(3) F；(4) F；(5) T。

(1)--(3) 是漸近符號（資結書第~1.5~節），(4) 是 AVL 樹（第~10~章），(5) 是優先佇列（第~8~章）。 本題放在佔 $6$ 分的第~1~章。

\textbf{(1) T}。 取 $c = 3$、$n_0 = 1$：$n \ge 1$ 時 $3n^3 + 7n^2 \ge 3n^3 \ge 3n^2$，符合 $\Omega(n^2)$ 的定義。 （上界下界都要的話，
它是 $\Theta(n^3)$；$\Omega(n^2)$ 只是一個比較鬆的下界，仍然成立。）

\textbf{(2) F}。 若 $n^3/\log^2 n = \Theta(n^3)$，就有常數 $c > 0$ 使得夠大的 $n$ 都有 $n^3/\log^2 n \ge c n^3$，即 $\log^2 n \le 1/c$，
但 $\log^2 n \to \infty$，矛盾。 正確的關係是 $n^3/\log^2 n = o(n^3)$，它是 $O(n^3)$ 但不是 $\Omega(n^3)$。

\textbf{(3) F}。 當 $f(n) \ge 1$（夠大的 $n$）時，$f(n) \le f(n)^2$，敘述成立；但題目對所有函數 $f$ 而言，反例是遞減到 $0$ 的函數。
取 $f(n) = 1/n$：$(f(n))^2 = 1/n^2$，若 $1/n \le c/n^2$ 對夠大的 $n$ 成立，就得 $n \le c$，矛盾。 所以一般而言不成立。 （這是 CLRS 的
經典是非題；若限定 $f$ 是執行時間這類 $\ge 1$ 的函數，答案才是 T。）

\textbf{(4) F}。 AVL 樹\textbf{插入}後最多一次（單或雙）旋轉就能恢復平衡，因為旋轉後那棵子樹的高度回到插入前的高度（資結書定理~10.2）。
\textbf{刪除}不一樣：旋轉後子樹的高度可能比刪除前少 $1$，不平衡會繼續往上傳，最壞要沿著路徑一路旋轉到根，需要 $\Theta(\log n)$ 次
（資結書第~10.2.6~節）。 例如高度 $h$ 的費氏樹（每個節點的左子樹都比右子樹高 $1$ 的最瘦 AVL 樹）刪掉適當的一個節點，要旋轉大約 $h/2$ 次
（$h = 11$、$232$ 個節點時最多 $5$ 次），次數隨 $h$ 成長，不是一次。

\textbf{(5) T}。 用排好序的串列做優先佇列，取出最大（或最小）值在串列頭，$O(1)$；但插入要先找到新元素的位置，最壞要走過整條串列，
$\Theta(n)$（資結書第~8.3~節）。 「linear time」指最壞情況，成立。 （若用陣列存排好序的資料，可以二分搜尋位置，但仍要挪動 $\Theta(n)$ 個
元素，一樣是線性。）
\end{pastsoln}
